% Baseline-fidelity response and corresponding manuscript text.
% Fill the result macros only from the frozen CARLA-bank analysis.  Do not
% replace them with trajectory-level impressions or unpaired summaries.
% Replace only this repository macro after the prepared local release tree has
% been committed and published. The current PC lacks the required `gh` CLI.
\providecommand{\BaselineCodeCommit}{\textbf{[INSERT IMMUTABLE PUBLIC COMMIT URL AFTER PUBLISHING]}}
\providecommand{\CarlaSceneCount}{five matched scene blocks (40 valid runs)}
\providecommand{\CarlaEmptyShadowFinding}{All retained controllers had 0/5
collisions and 0/5 near collisions in the empty-shadow condition.  DREAM's
mean-speed conservatism tax relative to IDEAM was 1.22~m\,s$^{-1}$ [1.13,
1.32].  Its maximum follower speed loss differed from IDEAM by only
$5.7\times10^{-6}$~m\,s$^{-1}$ [$-1.23\times10^{-3}$,
$1.25\times10^{-3}$], and the paired hard-braking actor count was unchanged.}
\providecommand{\CarlaThreatFinding}{In the true-threat condition, DREAM had
0/5 collisions and 2/5 near collisions, whereas IDEAM had 4/5 and 5/5.  The
ADA-sourced and APF-sourced controls remained collision-free but had 5/5 and
4/5 near collisions.  DREAM increased the mean episode-minimum signed
clearance by 1.19~m [0.89, 1.56], 0.67~m [0.52, 0.82], and 0.19~m [0.02,
0.39] relative to IDEAM, ADA-sourced, and APF-sourced, respectively.}
\providecommand{\CarlaFigureNumber}{\ref{fig:carla-two-condition-profiles}}
\providecommand{\BaselineFidelityTableNumber}{\ref{tab:baseline_fidelity}}

\noindent\textbf{Reviewer comment.}
\emph{Baseline fairness requires clarification. The baselines are
re-implemented on the authors' MPC--CBF backbone. In particular, OA-CMPC
appears to be stripped of its original dual-branch contingency structure,
which may affect its intended performance. The authors should document the
fidelity of each re-implementation, justify the modifications, and ideally
release the baseline code.}

\medskip
\noindent\textbf{Response.}
We agree with the reviewer.  Our audit confirmed that the submitted
OA-CMPC-labelled arm was not a faithful implementation of the published
OA-CMPC optimizer.  It retained only the tangent-line occlusion geometry and
time-to-collision-based reachable-set circles, rendered these circles as a
risk source $Q(\mathbf{x},t)$, and propagated that source through the DREAM
field and MPC--CBF stack.  It replaced the published dual-/multi-branch
optimizer by one worst-case branch.  Consequently, it did not implement
distinct nominal/exploration and contingency trajectories, branch-specific
decision variables, their shared-horizon consensus constraints, the ADMM
coordination updates, or the native OA-CMPC risk-boundary constraints.  These
omissions alter the method being compared and cannot be justified as a
performance-neutral implementation detail.

We therefore withdrew the OA-CMPC method-level comparison.  The revised
CARLA figure, aggregate statistics, and associated discussion contain no
OA-CMPC curve or ranking, and the earlier OA-CMPC panels and comparative
claims have been removed.  The archived adapter is now described only as an
``OA-inspired single-branch risk-source surrogate'' and is not used as
evidence about OA-CMPC.  A future native implementation would require the
published branch structure and solver; it is outside the evidential scope of
this revision.

We also revised the names and inferential scope of the retained comparison
arms.  IDEAM is the native no-field reference: it retains the IDEAM gap and
decision logic and its LMPC--CBF controller, while disabling the risk field
and all three DREAM-to-controller coupling channels.  The two other arms are
now named the \emph{ADA-sourced shared-backbone control} and the
\emph{APF-sourced shared-backbone control}.  They are controlled
source-substitution experiments, not end-to-end reproductions of the cited
ADA or APF planners.  In both arms, the CARLA observation adapter, IDEAM
decision/MPC--CBF backbone, field propagation equation, and the decision-veto,
MPC-cost, and CBF-modulation channels are held fixed.  Only the instantaneous
spatial source supplied to the propagation equation is changed.  The ADA
control uses the asymmetric momentum-wave source after per-class integral
matching and contains no explicit occlusion-shadow source.  The APF control
uses the repulsive potential after per-class integral matching; it omits an
attractive goal term because route generation and tracking remain common to
all arms, and it likewise contains no explicit occlusion-shadow source.  The
common propagation equation still transports and retains both substituted
sources over time.  All retained CARLA arms also share the same 10-Hz
trajectory tracker, lane-hold fallback, and visible-actor safety supervisor;
the comparison changes the high-level plan delivered to this common execution
layer.

Table~\BaselineFidelityTableNumber{} now records, for every retained arm and
the withdrawn OA surrogate, the method-specific element, common components,
adaptations, omissions, and valid interpretation.  The paired CARLA comparison
uses the same frozen nominal
scene manifest, stochastic seed, actor definitions, and traffic-control laws
for all retained arms.  Realized initial states must pass common predeclared
tolerances, and true-threat reveal is determined causally from semantic LiDAR
rather than imposed by a controller-specific schedule.  Controller order is
randomized within scene, and the scene seed is the sampling unit.  It comprises
\CarlaSceneCount{} across the empty-shadow and true-threat conditions.
Brackets denote descriptive 95\% complete-scene-block bootstrap intervals from
10,000 resamples; no population-level hypothesis test is claimed.  The
observed results are reported without extrapolation as follows:
\CarlaEmptyShadowFinding{}
\CarlaThreatFinding{} The revised code map, frozen manifests, runner, and
analysis scripts have been assembled in the local release tree.  The
immutable public URL to insert after publishing is \BaselineCodeCommit{}.

\medskip
\noindent\textbf{Corresponding manuscript edit (baseline definition).}

\begin{quote}
\textbf{Comparator scope and implementation fidelity.}
IDEAM is used as a native no-field reference and retains its gap-based
decision module and LMPC--CBF controller under the same vehicle model, route
request, and CARLA observation interface as DREAM.  The ADA-sourced and
APF-sourced arms are source-substitution controls rather than end-to-end
reproductions of the corresponding planners.  For these controls, field
transport and the three downstream couplings---decision veto, MPC field cost,
and risk-dependent CBF modulation---are retained, while only the spatial
source term is replaced.  The ADA source is an integral-matched asymmetric
momentum-wave field without an explicit shadow term.  The APF source is an
integral-matched repulsive potential without an explicit shadow term; its
attractive term is excluded because route generation and tracking are held
common.  Table~\ref{tab:baseline_fidelity} documents these adaptations and
their inferential scope.  An OA-inspired adapter was excluded because it maps
tangent-line reachable-set circles into a single risk source but does not
implement the published OA-CMPC dual-/multi-branch trajectories, consensus
constraints, ADMM coordination, or native risk-boundary constraints.
\end{quote}

\noindent\textbf{Corresponding manuscript edit (CARLA comparison).}

\begin{quote}
The matched CARLA experiment compares DREAM with the IDEAM no-field reference
and the ADA-sourced and APF-sourced controls.  Each controller receives the
same frozen nominal scene construction within a seed in both the empty-shadow
and true-threat conditions, and every realized initial state must satisfy the
same qualification gates.  Figure~\CarlaFigureNumber{} reports ego speed and
signed minimum oriented-box clearance relative to the semantic-LiDAR reveal
in each true-threat arm; its matching empty-shadow trace uses that
controller--seed reveal only as the common time origin.  Aggregate outcomes
treat the scene seed, rather than individual
simulation ticks, as the sampling unit.  \CarlaEmptyShadowFinding{}
\CarlaThreatFinding{}
\end{quote}
